% Fragmento público generado desde propietarios íntegros.
% Residencia: 52.9.2.
% Fuente: ../incoming/radion_monografia_20260827/manuscrito/sections/10_minkowski_hodge_lorentz.tex:101-186
\subsubsection{Réalisation lorentzienne des trajectoires hélicoïdales}

Dans un champ magnétique uniforme, une particule de charge \(q\), de masse \(m\) et de
facteur relativiste \(\gamma\) possède
\begin{equation}
\omega_c=\frac{|q|B}{\gamma m},
\qquad
r_L=\frac{p_\perp}{|q|B},
\qquad
\Delta z=\frac{2\pi p_\parallel}{|q|B}.
\label{eq:lorentz-helix}
\end{equation}
La phase cyclotronique se ferme tandis que le mouvement parallèle avance : l'orbite
est hélicoïdale. Le parallélisme avec \eqref{eq:H9-helix} est structurel et peut
être formalisé par un transducteur qui préserve phase, orientation et pas. On
n'identifie pas le quotient \(q_9\) à la coordonnée spatiale \(z\) ; on reconnaît
la même architecture de retour plus avance.

La force de Lorentz ne se dérive pas de l'image d'une spirale par ressemblance.
Elle se dérive en complétant \eqref{eq:lorentz-chain}. L'image aide à comprendre
pourquoi phase circulaire et mémoire longitudinale sont compatibles ; la forme de
champ détermine la dynamique.

\subsubsection{Confluence entre action, aire, information et gravité}

\paragraph{Réalisation Cartan--Holst de l'holonomie}

La promotion gravitationnelle s'exprime au moyen d'un morphisme d'holonomies
\begin{equation}
\operatorname{Hol}_{\mathcal A}
\bigl(\mathfrak R_{\mathrm{grav}}(\gamma)\bigr)
=\rho_C\bigl(\operatorname{Hol}_{\TPK}(\gamma)\bigr).
\label{eq:holonomy-gravity}
\end{equation}
Le membre de droite contient la mémoire de chemin ; \(\rho_C\) la réalise dans une
connexion gravitationnelle. La torsion de Cartan
\[
T^I=D_\omega e^I
\]
apparaît dans le codomaine. La torsion globale \(\Delta_4\) est une coordonnée
en amont qui peut alimenter l'application, mais n'est pas renommée \(T^I\) sans
\eqref{eq:holonomy-gravity}.

\paragraph{Transduction de l'aire d'action en aire informationnelle}

Le radion fixe un quantum d'action par phase. Le lecteur hexade--octade de
Barbero utilise ultérieurement les canaux \(q_\pm\), l'incidence \(90/120\) et
la fonctionnelle \eqref{eq:barbero-núcleo}. La direction causale est
\[
(A,C^*,\hret)\to q_\pm\to K_{6\to8}
\to\gamma_{\mathrm{BI}}\to\text{aire/information}.
\]
L'alphabet tritique apporte l'unité d'entropie \(\log3\) et une coupure de
\(81\) canaux apporte \(81\log3\). Cette entropie ne se confond ni avec
\(\log2\), ni avec \(\log\varphi\), ni avec une entropie microcanonique, ni avec une
entropie de von Neumann sans spécifier l'état.

La synthèse physique est affirmative mais typée : le radion apporte l'unité
d'action orientée qu'une holonomie peut transporter ; Barbero publie la
normalisation d'aire ; l'écran causal et Hodge permettent de réaliser des champs ;
Cartan--Holst reçoit l'holonomie en géométrie gravitationnelle.

\begin{statusbox}
L'équation scalaire du radion n'est pas à elle seule une équation d'Einstein. Sa
contribution à la gravité consiste à être une section exacte de l'architecture
d'action, d'orientation et d'holonomie que la réalisation gravitationnelle consomme.
\end{statusbox}

\subsubsection{Transformation de Wick comme changement de lecteur analytique}

Dans la formulation conventionnelle, la rotation de Wick relie une variable
temporelle lorentzienne à une coordonnée euclidienne par un prolongement
complexe. Dans HMT, la racine \(J^2=-I\), l'opérateur involutif \(R^2=I\) et le seuil
\(N^2=0\) existent déjà auparavant. Le changement de lecteur
\[
J\longleftrightarrow R
\]
transporte un couple commun/orienté entre publication elliptique et
publication hyperbolique. Le prolongement complexe reconnaît ensuite cette
opération en coordonnées analytiques.

Cela explique pourquoi cercle, ellipse d'action et intérieur hyperbolique peuvent
appartenir au même état sans être la même métrique. Le cercle utilise \(J\) pour
la phase ; la forme de l'ellipse utilise une compression symplectique générée par \(R\) ; Klein mesure
l'intérieur par le birapport ; l'écran de Minkowski réalise la causalité
dans un autre codomaine.
